\begin{abstract}
LLM-as-judge is increasingly used to evaluate code correctness without test execution, but practitioners lack guidance on how model scale affects reliability. We present the first scale-controlled analysis of error patterns in LLM code judges, comparing 7B, 70B, and proprietary models under fixed evaluation settings. Our key finding is that scale predicts error \textit{type}, not just error rate: smaller models systematically over-accept (FPR=74.8\%), while larger models systematically under-accept (FNR=29.3\%). This FPR-FNR tradeoff is statistically significant ($\chi^2$=45.78, $p$=3.27$\times$10$^{-8}$) and explains two surprising results. First, ensemble voting across scales \textit{degrades} accuracy by 8.05\% compared to the best single judge (McNemar $p$=1.45$\times$10$^{-6}$)---the over-accepting majority drowns the accurate minority signal. Second, unanimous agreement correlates with \textit{lower} accuracy (35.8\%) than disagreement (39.7\%), because unanimous ``correct'' verdicts capture shared over-acceptance bias. We recommend using 70B models for optimal cost-accuracy tradeoff (capturing 95\% of scale benefit) and selecting scale based on error-cost asymmetry rather than naive ensemble combination.
\end{abstract}
